\chapter{Peanut allergy}
\index{Peanut allergy}

\medskip
\noindent\textit{Educational information; seek professional advice for individual care.}

\section*{What it is}
Peanut allergy is an immune reaction to peanut proteins. Reactions can be mild or life-threatening and may change from one exposure to another. Diagnosis should be based on the history and appropriate allergy assessment rather than a screening test alone.

\section*{Possible reactions}
Symptoms may include itchy mouth, hives, swelling, vomiting, cough, wheeze, hoarse voice, throat tightness, dizziness, fainting, or breathing difficulty. Anaphylaxis can occur without skin symptoms and can progress quickly.

\section*{Assessment and prevention}
An allergy clinician may use the exposure history, examination, skin or blood testing, and sometimes a supervised food challenge. Avoid peanut only as advised, read ingredient information, consider cross-contact, and tell schools, workplaces, carers, and travel companions. Keep a written emergency plan and prescribed adrenaline or epinephrine device available and learn how to use it.

\section*{When to Seek Medical Care}
For a suspected severe reaction, use the prescribed adrenaline or epinephrine immediately according to the person’s action plan and contact local emergency services. Lay the person down if possible, or allow a person who is struggling to breathe to sit with legs extended; do not let them stand or walk. A second prescribed dose may be needed if symptoms persist or return. Antihistamines do not replace adrenaline for anaphylaxis.

\section*{Follow-up}
After any significant reaction, obtain medical assessment even if symptoms improve. Discuss ongoing avoidance, emergency training, and whether supervised oral immunotherapy is appropriate; do not try it at home.

\section*{Sources}
\begin{itemize}
\item NCBI Bookshelf, Peanut Allergy: \url{https://www.ncbi.nlm.nih.gov/books/NBK538526/}
\item Food Allergy Research and Education reaction guidance: \url{https://www.foodallergy.org/resources/recognizing-and-treating-reaction-symptoms}
\end{itemize}
